\section{More on metric spaces: continuity}
A metric $d$ measures distance. The open ball $B_r(s)$ consists of the
points whose distance from $s$ is less than $r$. A continuous function
sends points sufficiently close to a \emph{fixed} point to points close
to its image. Uniform continuity uses one input tolerance for all pairs
of points. An open set in a subspace is understood relative to that
subspace; it need not be open in the surrounding Euclidean space.

\exercise{21.1}{Uniform continuity of a vector-valued function}
\goal Show that $\gamma(t)=(f_1(t),\ldots,f_k(t))$ is uniformly
continuous when all its coordinate functions are uniformly continuous.
\idea Make each coordinate error small enough that the sum of their
squares is less than the desired squared distance.
\step{Choose the input tolerance before choosing the points.}
Fix $\varepsilon>0$. For each $j$, uniform continuity of $f_j$ gives
$\delta_j>0$ such that, for \emph{all} $s,t\in[a,b]$,
\[
 |s-t|<\delta_j\ \Longrightarrow\
 |f_j(s)-f_j(t)|<\varepsilon/\sqrt{k}.
\]
The minimum $\delta=\min(\delta_1,\ldots,\delta_k)$ is positive because
there are only finitely many coordinates. If $|s-t|<\delta$, every
coordinate estimate holds. Consequently
\[
 d(\gamma(s),\gamma(t))
 =\left(\sum_{j=1}^k|f_j(s)-f_j(t)|^2\right)^{1/2}
 <\left(k\frac{\varepsilon^2}{k}\right)^{1/2}=\varepsilon.
\]
\answer{The same $\delta$ works for every pair $s,t$, which is precisely
uniform continuity.}

\exercise{21.2}{Continuity at one point in terms of open sets}
\goal Prove both directions of the open-neighborhood characterization
at $s_0\in S$.
\step{Continuity implies the neighborhood condition.}
Suppose $f$ is continuous at $s_0$ and $U\subseteq S^*$ is open with
$f(s_0)\in U$. Openness provides $\varepsilon>0$ for which
$B^*_{\varepsilon}(f(s_0))\subseteq U$. Continuity supplies $\delta>0$
with $d(s,s_0)<\delta$ implying
$d^*(f(s),f(s_0))<\varepsilon$. The open set $V=B_\delta(s_0)$ contains
$s_0$ and satisfies $f(V)\subseteq U$.
\step{The neighborhood condition implies continuity.}
Conversely, fix $\varepsilon>0$ and apply the assumed condition to the
open ball $U=B^*_{\varepsilon}(f(s_0))$. We obtain an open $V$ containing
$s_0$ with $f(V)\subseteq U$. Since $V$ is open, it contains a ball
$B_\delta(s_0)$ for some $\delta>0$. Hence
\[
 d(s,s_0)<\delta\ \Longrightarrow\ s\in V
 \ \Longrightarrow\ f(s)\in U
 \ \Longrightarrow\ d^*(f(s),f(s_0))<\varepsilon.
\]
\answer{This is exactly the metric definition of continuity at $s_0$.}

\exercise{21.3}{Distance from a fixed point is uniformly continuous}
\goal For fixed $s_0$, prove uniform continuity of $f(s)=d(s,s_0)$.
\step{Derive the reverse triangle inequality.}
The triangle inequality gives
$d(s,s_0)\le d(s,t)+d(t,s_0)$, so
$d(s,s_0)-d(t,s_0)\le d(s,t)$. Interchanging $s,t$ bounds the negative
of the same difference. Together these statements give
\[
 |f(s)-f(t)|=|d(s,s_0)-d(t,s_0)|\le d(s,t).
\]
Given $\varepsilon>0$, choose $\delta=\varepsilon$. Whenever
$d(s,t)<\delta$, the displayed inequality implies
$|f(s)-f(t)|<\varepsilon$.
\answer{Distance from a fixed point is even Lipschitz with constant $1$;
this means its output distance never exceeds its input distance.}

\exercise{21.4}{Rational intervals suffice to test continuity}
\goal Establish the equivalence of continuity, openness of preimages
of all open intervals, and openness for intervals with rational endpoints.
\step{The first two implications.}
If $f$ is continuous, Theorem~21.3 makes the inverse image of every
open subset of $\R$ open in $S$. In particular, this holds for every
interval $(a,b)$. Restricting $a,b$ to rationals gives the third condition.
\step{Recover continuity using only rational endpoints.}
Assume $f^{-1}((a,b))$ is open whenever $a<b$ are rational. Fix
$s_0\in S$ and $\varepsilon>0$. Density of $\Q$ gives rationals with
\[
 f(s_0)-\varepsilon<a<f(s_0)<b<f(s_0)+\varepsilon.
\]
The set $V=f^{-1}((a,b))$ is open and contains $s_0$, so some
$B_\delta(s_0)$ is contained in $V$. If $d(s,s_0)<\delta$, then
$a<f(s)<b$, and therefore $|f(s)-f(s_0)|<\varepsilon$.
\answer{Continuity at every $s_0$ follows. A dense set of possible
interval endpoints gives arbitrarily small neighborhoods of every output.}

\exercise{21.5}{What can fail on a noncompact Euclidean set}
\goal Construct (a) an unbounded continuous function and (b) a bounded
continuous function with no maximum on a noncompact $E\subseteq\R^k$.
\idea By Heine--Borel, a compact Euclidean set is exactly a closed
bounded set. Failure of compactness gives either points escaping to
infinity or a missing limit point.
\step{(a) Construct an unbounded function.}
If $E$ is unbounded, set $h(x)=\|x\|$, the Euclidean distance from
zero. Exercise~21.3 proves continuity. For each $M>0$, unboundedness
of $E$ supplies $x\in E$ with $h(x)>M$.

If $E$ is bounded, it cannot be closed. Choose
$x_0\in\overline E\setminus E$ and set $h(x)=1/\|x-x_0\|$ on $E$.
The denominator is positive on $E$, and the distance function and the
reciprocal function on $(0,\infty)$ are continuous, so $h$ is continuous.
Since $x_0\in\overline E$, for every $M>0$ some $x\in E$ satisfies
$0<\|x-x_0\|<1/M$; thus $h(x)>M$.
\step{(b) Compress the unbounded values into a bounded interval.}
In either case $h\ge0$ and $h$ is unbounded above. Define
\[
 g(x)=\frac{h(x)}{1+h(x)}=1-\frac1{1+h(x)}.
\]
This is continuous and $0\le g(x)<1$. For any $x\in E$, choose $y\in E$
with $h(y)>h(x)$. Since $t/(1+t)$ is strictly increasing for $t\ge0$
(as follows by cross multiplication), $g(y)>g(x)$. Thus no point
can maximize $g$.
\answer{Compactness is essential both for boundedness of continuous
functions and for attainment of their extreme values.}

\exercise{21.6}{Compositions of continuous metric-space functions}
\goal Show that $g\circ f:S_1\to S_3$ is continuous when both component
functions are continuous.
\step{Pull an open set back in two stages.}
Let $U\subseteq S_3$ be open. Continuity of $g$ says $g^{-1}(U)$ is
open in $S_2$. Continuity of $f$ then says
$f^{-1}(g^{-1}(U))$ is open in $S_1$. The identity
\[
 (g\circ f)^{-1}(U)=f^{-1}(g^{-1}(U))
\]
holds because each side consists of exactly those $s$ with $g(f(s))\in U$.
Theorem~21.3 therefore proves continuity of $g\circ f$.
\answer{The order of inverse images is the reverse of the order in
which the functions are applied.}

\exercise{21.7}{Restricted metrics and continuity of a path}
\step{(a) Restricting a metric.}
For $s,t,u\in E\subseteq S$, the numbers $d(s,t)$ are nonnegative,
$d(s,t)=0$ exactly when $s=t$, symmetry holds, and
$d(s,u)\le d(s,t)+d(t,u)$. These statements already hold in $S$,
so they still hold when the available points are restricted to $E$.
Thus $(E,d)$ is a metric space. In particular $d(x,y)=|x-y|$ is a
metric on every subset of $\R$.
\step{(b) Spell out the definition for a path.}
A function $\gamma:[a,b]\to\R^k$ is continuous if, for every
$t_0\in[a,b]$ and every $\varepsilon>0$, there is $\delta>0$ such that
\[
 t\in[a,b],\quad |t-t_0|<\delta
 \quad\Longrightarrow\quad
 \left(\sum_{j=1}^k|\gamma_j(t)-\gamma_j(t_0)|^2\right)^{1/2}
 <\varepsilon.
\]
\answer{At an endpoint, only points of $[a,b]$ are tested. The definition
does not ask for values of the path outside its domain.}

\exercise{21.8}{Uniform continuity sends Cauchy sequences to Cauchy sequences}
\goal Prove that $(f(s_n))$ is Cauchy in $S^*$ when $(s_n)$ is Cauchy
in $S$ and $f$ is uniformly continuous.
\step{Match the two definitions.}
Fix $\varepsilon>0$. Uniform continuity gives $\delta>0$ such that
$d(s,t)<\delta$ implies $d^*(f(s),f(t))<\varepsilon$, for all $s,t\in S$.
Since $(s_n)$ is Cauchy, some $N$ has $d(s_m,s_n)<\delta$ whenever
$m,n>N$. Consequently
\[
 m,n>N\ \Longrightarrow\ d^*(f(s_m),f(s_n))<\varepsilon.
\]
\answer{This is the Cauchy condition in the target. We did not assume
either sequence has a limit. Mere continuity would not suffice: $1/x$
on $(0,1)$ sends the Cauchy sequence $1/(n+1)$ to the non-Cauchy
sequence $n+1$.}

\exercise{21.9}{Continuous surjections involving a square}
\step{(a) Project the square onto an interval.}
On $[0,1]^2$, put $p(x_1,x_2)=x_1$. The inequality
$|p(x)-p(y)|\le\|x-y\|$ proves continuity. Every $u\in[0,1]$ is the
image of $(u,0)$, so the range is exactly $[0,1]$.
\step{(b) Can an interval fill a square?}
Yes: a continuous surjection $[0,1]\to[0,1]^2$ is called a
\emph{space-filling curve}. Dimension alone does not prevent one.
The question asks for a judgment, rather than a formula; here is an
existence argument explaining why the answer is not merely a guess.

Let $C$ be the middle-thirds Cantor set. Each point of $C$ has a
unique expansion $t=\sum_{j\ge1}2b_j3^{-j}$ with $b_j\in\{0,1\}$.
Define
\[
 H(t)=\left(\sum_{j\ge1}b_{2j-1}2^{-j},\
                 \sum_{j\ge1}b_{2j}2^{-j}\right).
\]
Any two numbers in $[0,1]$ have binary expansions; interlacing their
digits produces a point of $C$ with precisely those two images.
Thus $H(C)=[0,1]^2$. If two Cantor points have their first $2m$
ternary digits equal, their coordinate differences are at most
$\sum_{j>m}2^{-j}=2^{-m}$. Conversely, points in distinct cylinders
specified by the first $2m$ digits are separated by a positive gap.
These finite-cylinder gaps prove continuity of $H$ on $C$.

Extend $H$ to each removed interval $(u,v)$ by straight-line
interpolation between $H(u)$ and $H(v)$. The resulting map stays in
the square because line segments between square points stay in the
square. It is continuous inside each gap. To check continuity at
$t\in C$, consider any sequence $t_n\to t$. For points in $C$ use
continuity of $H$. For points in gaps, either a subsequence lies in a
fixed gap whose endpoint is $t$ (where linear interpolation is
continuous), or the gap lengths tend to zero. In the latter case both
endpoints tend to $t$, and their images tend to $H(t)$, so every
convex combination does too. Only finitely many gaps have length at
least any prescribed positive number; hence these cases cover every
possible subsequence. The sequential criterion proves continuity.
\answer{The extension is a continuous map on $[0,1]$ and is already
onto the square on its Cantor subset. A curve can therefore fill a square.}

\exercise{21.10}{Four questions about continuous ranges}
\step{(a) Map $(0,1)$ onto $[0,1]$.}
Use $f(t)=(1+\sin(4\pi t))/2$. This continuous function takes values
only in $[0,1]$, attains $1$ at $t=1/8$, and attains $0$ at $t=3/8$.
The intermediate value theorem on $[1/8,3/8]$ supplies every value
between them. Hence the map is onto.
\step{(b) Map $(0,1)$ onto $\R$.}
Use $g(t)=\tan(\pi(t-1/2))$. Its argument lies in
$(-\pi/2,\pi/2)$, where tangent is continuous and strictly increasing,
with limits $-\infty$ and $+\infty$ at the two ends. For any real
$y$, choose interior points where $g$ is below and above $y$ and apply
the intermediate value theorem.
\step{(c) Map $[0,1]\cup[2,3]$ onto $[0,1]$.}
Define $h(t)=t$ on $[0,1]$ and $h(t)=t-2$ on $[2,3]$. On each
component this is continuous. A ball of radius less than $1$ around
any domain point cannot reach the other component, so continuity on
each component proves continuity on the whole domain. Each component
already maps onto $[0,1]$.
\step{(d) Rule out two proposed ranges.}
A continuous image of the compact interval $[0,1]$ is compact.
Neither $(0,1)$ (not closed in $\R$) nor $\R$ (unbounded) is compact.
If a continuous map were onto either, its image would be both compact
and noncompact, a contradiction.
\answer{The formulas supply (a)--(c); preservation of compactness
excludes both maps in (d).}

\exercise{21.11}{The rationals are not a countable intersection of open sets}
\goal Use the Baire theorem to rule out $\Q=\bigcap_{n\ge1}U_n$ with
all $U_n$ open in $\R$.
\idea Such $U_n$ would all be dense. Intersect them also with the
complements of individual rationals and reach an empty dense set.
\step{Construct dense open sets.}
Suppose the representation exists. Since every $U_n$ contains $\Q$
and $\Q$ is dense, $U_n$ is dense. List the rationals as
$\Q=\{q_1,q_2,\ldots\}$. Each $V_n=\R\setminus\{q_n\}$ is open
and dense: every nonempty open interval contains points other than a
single fixed point.
\step{Apply completeness exactly where needed.}
The real line is a complete metric space. Theorem~21.7(a), applied
to the interlaced sequence $U_1,V_1,U_2,V_2,\ldots$, says their
intersection is dense. But it equals
\[
 \left(\bigcap_nU_n\right)\cap\left(\bigcap_nV_n\right)
 =\Q\cap(\R\setminus\Q)=\varnothing,
\]
which is not dense in the nonempty space $\R$.
\step{Take complements for the equivalent formulation.}
If $\R\setminus\Q=\bigcup_nF_n$ with each $F_n$ closed, De Morgan's
law would give $\Q=\bigcap_n(\R\setminus F_n)$ with open terms.
\answer{Neither proposed representation is possible.}

\exercise{21.12}{Infinitely many disjoint sets of second category}
\goal Give pairwise disjoint subsets of $\R$, each not a countable union
of nowhere dense subsets of $\R$.
\step{Choose intervals and invoke the right Baire consequence.}
For $n\ge1$, let $E_n=(2n,2n+1)$. These intervals are pairwise
disjoint. Suppose some $E_n$ were a countable union of nowhere dense
sets. By Theorem~21.7(c), its complement would be dense in $\R$.
But that complement misses the nonempty open interval $E_n$,
contradicting density.
\answer{The sequence $((2n,2n+1))_{n\ge1}$ works. The category here
is category in the ambient space $\R$, as the exercise requires.}

\exercise{21.13}{Oscillation detects continuity}
\goal For
\[
 \omega_f(x)=\inf_{\delta>0}\sup\{|f(y)-f(z)|:
                        y,z\in(x-\delta,x+\delta)\},
\]
prove $f$ is continuous at $x$ exactly when $\omega_f(x)=0$.
The suprema may be infinite when $f$ is locally unbounded; all
quantities are nonnegative.
\step{Continuity makes the oscillation zero.}
Fix $\varepsilon>0$. Continuity gives $\delta>0$ such that
$|f(y)-f(x)|<\varepsilon/4$ for $|y-x|<\delta$. For two such
points, the triangle inequality gives
$|f(y)-f(z)|<\varepsilon/2$. Taking the supremum can change a strict
inequality into a nonstrict one, but still gives a bound at most
$\varepsilon/2<\varepsilon$. Therefore
$0\le\omega_f(x)<\varepsilon$. Since this holds for every
$\varepsilon>0$, $\omega_f(x)=0$.
\step{Zero oscillation gives continuity.}
If the nonnegative infimum is zero, then for every $\varepsilon>0$
some $\delta>0$ has the corresponding supremum less than
$\varepsilon$. Set $z=x$, which is permitted because $x$ belongs to
the interval. Then
\[
 |y-x|<\delta\ \Longrightarrow\
 |f(y)-f(x)|\le\sup_{u,v\in(x-\delta,x+\delta)}|f(u)-f(v)|
 <\varepsilon.
\]
\answer{This is continuity at $x$. It is important that the interval
in the oscillation definition includes its center.}

\exercise{21.14}{Reduce every interval to the bounded closed case}
\goal Complete Theorem~21.11 after its proof for $[a,b]$.
\step{Assume a forbidden decomposition of an arbitrary interval.}
Suppose a nondegenerate interval $I$ is a disjoint union of at least
two nondegenerate closed intervals $J_\alpha$. Choose two distinct
members. Select an interior point $x$ of one and an interior point
$y$ of the other; relabel so $x<y$. Because $I$ is an interval,
$[x,y]\subseteq I$.
\step{Intersect every piece with $[x,y]$.}
Discard empty intersections. The remaining
$K_\alpha=J_\alpha\cap[x,y]$ are pairwise disjoint closed intervals
whose union is $[x,y]$. The interval containing $x$ has positive
length to the right of $x$, and the interval containing $y$ has
positive length to the left of $y$, because $x,y$ were chosen in
interiors. Every other surviving interval lies strictly between
$x$ and $y$ and retains its original positive length. Thus none of
the surviving intersections is a singleton. There are at least two
of them, since the intervals containing $x$ and $y$ are distinct.
\step{Use the case already established.}
We have written the bounded closed interval $[x,y]$ as a disjoint
union of at least two nondegenerate closed intervals. This contradicts
the bounded closed case of Theorem~21.11.
\answer{The theorem holds for all nondegenerate intervals, including
open, half-open, and unbounded intervals. The interior-point choice
prevents unwanted singleton pieces at the truncation endpoints.}
